% Section 03 - Génération de relevés de notes
\section{Génération de relevés de notes}

\subsection{Vue d'ensemble}

Le menu \menu{Générer les relevés} permet de créer automatiquement des relevés de notes personnalisés pour vos étudiants à partir d'un fichier Excel contenant les notes.

\textbf{Prérequis :}
\begin{itemize}[leftmargin=*]
    \item Avoir configuré au moins une classe dans \menu{Config des relevés}
    \item Disposer d'un fichier Excel avec les données des étudiants et leurs notes
    \item Avoir configuré les entêtes de l'établissement
\end{itemize}

\subsection{Préparation du fichier Excel}

\subsubsection{Structure obligatoire}

Votre fichier Excel doit contenir au minimum les colonnes suivantes :

\begin{tcolorbox}[colback=blue!5!white,colframe=blue!75!black,title=Colonnes obligatoires]
\begin{longtable}{|l|p{10cm}|}
\hline
\textbf{Colonne} & \textbf{Description} \\
\hline
\endfirsthead
\hline
\textbf{Colonne} & \textbf{Description} \\
\hline
\endhead
\hline
\endfoot

\texttt{MATRICULE} & Numéro matricule unique de l'étudiant \\
\hline
\texttt{NOM} & Nom de famille (sera affiché en MAJUSCULES) \\
\hline
\texttt{PRENOM} & Prénom(s) de l'étudiant \\
\hline
\texttt{SEMESTRE} & Code du semestre (ex: S1, S2, S3, S4, S5, S6) \\
\hline
\texttt{NIVEAU} & Niveau d'études (ex: L1, L2, L3, M1, M2) \\
\hline
\texttt{FILIERE} & Code ou intitulé de la filière \\
\hline
\end{longtable}
\end{tcolorbox}

\subsubsection{Colonnes de notes}

En plus des colonnes obligatoires, ajoutez une colonne pour chaque EC (Élément Constitutif) défini dans votre configuration de classe.

\textbf{Exemple :} Si votre configuration contient les ECs suivants :
\begin{itemize}
    \item Mathématiques
    \item Physique
    \item Chimie
    \item Biologie
\end{itemize}

Votre Excel devra contenir les colonnes :
\begin{verbatim}
MATRICULE | NOM | PRENOM | SEMESTRE | NIVEAU | FILIERE | Mathématiques |
Physique | Chimie | Biologie
\end{verbatim}

\begin{figure}[h]
    \centering
    % \includegraphics[width=\textwidth]{images/excel_structure_releves.png}
    \fbox{\parbox{0.95\textwidth}{\centering [CAPTURE D'ÉCRAN : EXEMPLE DE FICHIER EXCEL POUR RELEVÉS]\\ \vspace{0.3cm} Montrant les colonnes obligatoires + colonnes de notes}}
    \caption{Structure d'un fichier Excel pour génération de relevés}
\end{figure}

\subsubsection{Colonnes optionnelles}

Vous pouvez ajouter des colonnes supplémentaires qui ne seront pas utilisées par l'application :

\begin{itemize}[leftmargin=*]
    \item \texttt{DATE\_NAISSANCE} : Date de naissance
    \item \texttt{LIEU\_NAISSANCE} : Lieu de naissance
    \item \texttt{SEXE} : M/F
    \item \texttt{EMAIL} : Adresse email
    \item Toute autre colonne pour votre usage interne
\end{itemize}

\attention{Nommez vos colonnes de notes EXACTEMENT comme les noms d'ECs dans votre configuration de classe (sensible à la casse).}

\subsubsection{Format des notes}

\begin{itemize}[leftmargin=*]
    \item Les notes doivent être des nombres (format numérique Excel)
    \item Les notes peuvent être décimales : \texttt{15.5}, \texttt{12.75}
    \item Pour une note manquante, laissez la cellule vide ou mettez \texttt{ABS} (absent)
    \item Les notes doivent correspondre à la base définie dans la configuration (sur 10, 20, ou 100)
\end{itemize}

\subsection{Workflow complet de génération}

\subsubsection{Étape 1 : Accéder au menu}

\begin{enumerate}[leftmargin=*]
    \item Cliquez sur \menu{Générer les relevés} dans la barre latérale
    \item L'interface de génération s'affiche avec plusieurs onglets
\end{enumerate}

\subsubsection{Étape 2 : Charger le fichier Excel}

\begin{enumerate}[leftmargin=*]
    \item Dans l'onglet \textbf{Générateur}, cliquez sur \bouton{Charger fichier Excel}
    \item Sélectionnez votre fichier .xlsx ou .xls
    \item L'application valide automatiquement :
    \begin{itemize}
        \item Présence des colonnes obligatoires
        \item Format des données
        \item Nombre d'étudiants détectés
    \end{itemize}
    \item Un message de confirmation s'affiche avec le nombre de lignes importées
\end{enumerate}

\begin{figure}[h]
    \centering
    % \includegraphics[width=0.9\textwidth]{images/import_excel_releves.png}
    \fbox{\parbox{0.85\textwidth}{\centering [CAPTURE D'ÉCRAN : IMPORT EXCEL AVEC VALIDATION]\\ \vspace{0.3cm} Message : "✓ 150 étudiants importés avec succès"}}
    \caption{Import et validation du fichier Excel}
\end{figure}

\subsubsection{Étape 3 : Sélectionner la configuration de classe}

\begin{enumerate}[leftmargin=*]
    \item Dans le champ \champ{Sélectionner une configuration}, choisissez la classe correspondante
    \item Exemple : "Licence 1 Médecine - 2023-2024"
    \item La liste des semestres disponibles s'affiche automatiquement
\end{enumerate}

\info{Si vous n'avez pas encore créé de configuration, cliquez sur \bouton{Créer une nouvelle configuration} pour être redirigé vers le menu \menu{Config des relevés}.}

\subsubsection{Étape 4 : Choisir le(s) semestre(s)}

Vous avez deux options :

\paragraph{Option A : Un seul semestre}

\begin{enumerate}[leftmargin=*]
    \item Dans la liste déroulante \champ{Sélectionner un semestre}, choisissez le semestre souhaité
    \item Exemple : "Semestre 1 (S1)"
    \item Seuls les étudiants de ce semestre seront traités
\end{enumerate}

\paragraph{Option B : Tous les semestres (génération multi-semestres)}

\begin{enumerate}[leftmargin=*]
    \item Sélectionnez \textbf{"Tous les semestres"} dans la liste déroulante
    \item Une interface de sélection multi-semestres apparaît
    \item Cochez les semestres à inclure dans la génération
    \item Utilisez \bouton{Tout sélectionner} ou \bouton{Tout désélectionner} pour faciliter la sélection
\end{enumerate}

\begin{figure}[h]
    \centering
    % \includegraphics[width=\textwidth]{images/selection_multi_semestres.png}
    \fbox{\parbox{0.9\textwidth}{\centering [CAPTURE D'ÉCRAN : INTERFACE MULTI-SEMESTRES]\\ \vspace{0.3cm} Cases à cocher pour S1, S2, S3, S4, S5, S6\\ Boutons "Tout sélectionner" et "Tout désélectionner"}}
    \caption{Sélection de plusieurs semestres simultanément}
\end{figure}

\subsubsection{Étape 5 : Mapper les colonnes Excel aux ECs}

Si c'est la première fois que vous utilisez ce fichier Excel avec cette configuration, vous devez associer chaque EC à sa colonne Excel correspondante.

\begin{enumerate}[leftmargin=*]
    \item Cliquez sur l'onglet \textbf{Mapping des colonnes}
    \item Pour chaque EC listé, sélectionnez la colonne Excel correspondante dans la liste déroulante
    \item L'application propose automatiquement les colonnes dont le nom correspond
    \item Si un EC a une session de rattrapage :
    \begin{itemize}
        \item Activez le toggle \champ{Session}
        \item Mappez également la colonne de la session
    \end{itemize}
    \item Cliquez sur \bouton{Enregistrer le mapping}
\end{enumerate}

\begin{figure}[h]
    \centering
    % \includegraphics[width=\textwidth]{images/mapping_colonnes.png}
    \fbox{\parbox{0.9\textwidth}{\centering [CAPTURE D'ÉCRAN : INTERFACE DE MAPPING]\\ \vspace{0.3cm}
    Tableau avec colonnes : EC | Colonne Excel | Session | Colonne Session\\
    Mathématiques → [Dropdown: Mathématiques]\\
    Physique → [Dropdown: Physique Générale]}}
    \caption{Mapping des colonnes Excel vers les ECs de la configuration}
\end{figure}

\astuce{Le mapping est mémorisé ! La prochaine fois que vous chargerez un fichier Excel avec les mêmes noms de colonnes, le mapping se fera automatiquement.}

\subsubsection{Étape 6 : Configurer les options de génération}

Dans l'onglet \textbf{Options}, configurez les paramètres suivants :

\paragraph{Format d'exportation}

Choisissez parmi les 3 formats disponibles :

\begin{description}[leftmargin=3.5cm]
    \item[ZIP (défaut)] Archive compressée contenant tous les PDFs individuels
    \item[Fichiers individuels] Chaque relevé téléchargé séparément (sans archive)
    \item[PDF unique] Tous les relevés fusionnés en un seul fichier PDF
\end{description}

\paragraph{Compression}

\begin{itemize}[leftmargin=*]
    \item Activez \champ{Utiliser la compression} pour réduire la taille de l'archive ZIP
    \item Gain typique : 30-50\% de réduction de taille
    \item Légère augmentation du temps de génération
\end{itemize}

\paragraph{Format de nommage}

Choisissez comment les fichiers PDF seront nommés :

\begin{description}[leftmargin=3.5cm]
    \item[Détaillé (défaut)] \texttt{Releve\_NOM\_Prenom\_S1\_2023-2024.pdf}
    \item[Simple] \texttt{Releve\_MATRICULE.pdf}
\end{description}

\paragraph{QR Code chiffré (optionnel)}

\begin{itemize}[leftmargin=*]
    \item Activez \champ{Ajouter un QR code chiffré} pour intégrer un QR code sécurisé
    \item Le QR contient des informations chiffrées avec AES-128-ECB
    \item Clé de chiffrement basée sur le matricule de l'étudiant
    \item Permet la vérification d'authenticité du relevé
\end{itemize}

\begin{figure}[h]
    \centering
    % \includegraphics[width=0.8\textwidth]{images/options_generation_releves.png}
    \fbox{\parbox{0.75\textwidth}{\centering [CAPTURE D'ÉCRAN : ONGLET OPTIONS]\\ \vspace{0.3cm}
    Radio buttons : ○ ZIP ● Fichiers individuels ○ PDF unique\\
    ☑ Utiliser la compression\\
    Format nommage : [Dropdown: Détaillé]\\
    ☐ Ajouter un QR code chiffré}}
    \caption{Configuration des options de génération}
\end{figure}

\subsubsection{Étape 7 : Prévisualiser (optionnel)}

Avant de générer tous les relevés :

\begin{enumerate}[leftmargin=*]
    \item Cliquez sur l'onglet \textbf{Aperçu}
    \item Sélectionnez un étudiant dans la liste
    \item Cliquez sur \bouton{Générer l'aperçu}
    \item Un PDF de prévisualisation s'affiche dans le navigateur
    \item Vérifiez :
    \begin{itemize}
        \item Les informations de l'étudiant (nom, prénom, matricule)
        \item Les notes affichées pour chaque EC
        \item Le calcul des moyennes par UE
        \item La moyenne générale et la mention
        \item La mise en page et le thème
    \end{itemize}
    \item Si tout est correct, revenez à l'onglet \textbf{Générateur}
\end{enumerate}

\begin{figure}[h]
    \centering
    % \includegraphics[width=0.7\textwidth]{images/apercu_releve.png}
    \fbox{\parbox{0.65\textwidth}{\centering [CAPTURE D'ÉCRAN : APERÇU D'UN RELEVÉ PDF]\\ \vspace{0.3cm}
    Affichage d'un relevé complet avec notes, moyennes, mention}}
    \caption{Aperçu d'un relevé avant génération massive}
\end{figure}

\subsubsection{Étape 8 : Générer les relevés}

\begin{enumerate}[leftmargin=*]
    \item Retournez sur l'onglet \textbf{Générateur}
    \item Vérifiez que toutes les configurations sont correctes :
    \begin{itemize}
        \item Fichier Excel chargé ✓
        \item Configuration sélectionnée ✓
        \item Semestre(s) choisi(s) ✓
        \item Mapping des colonnes effectué ✓
        \item Options configurées ✓
    \end{itemize}
    \item Cliquez sur le bouton \bouton{Générer les relevés}
    \item Une barre de progression s'affiche montrant :
    \begin{itemize}
        \item Nombre de relevés générés / total
        \item Pourcentage de progression
        \item Temps estimé restant
    \end{itemize}
    \item Attendez la fin du traitement (temps variable selon le nombre d'étudiants)
\end{enumerate}

\begin{figure}[h]
    \centering
    % \includegraphics[width=0.8\textwidth]{images/progression_generation.png}
    \fbox{\parbox{0.75\textwidth}{\centering [CAPTURE D'ÉCRAN : BARRE DE PROGRESSION]\\ \vspace{0.3cm}
    Génération en cours : 47/150 relevés (31\%)\\
    ▓▓▓▓▓▓▓░░░░░░░░░░░░░░\\
    Temps estimé : 2 min 15 s}}
    \caption{Barre de progression pendant la génération}
\end{figure}

\subsubsection{Étape 9 : Téléchargement}

Une fois la génération terminée :

\begin{itemize}[leftmargin=*]
    \item \textbf{Format ZIP} : Une archive est automatiquement téléchargée
    \begin{itemize}
        \item Nom du fichier : \texttt{Releves\_[Date]\_[Heure].zip}
        \item Ouvrez l'archive et extrayez les PDFs
    \end{itemize}

    \item \textbf{Format Fichiers individuels} : Les fichiers sont téléchargés un par un
    \begin{itemize}
        \item Un délai de 100 ms entre chaque téléchargement évite de saturer le navigateur
        \item Les fichiers apparaissent dans votre dossier de téléchargements
    \end{itemize}

    \item \textbf{Format PDF unique} : Un seul fichier fusionné est téléchargé
    \begin{itemize}
        \item Nom : \texttt{Releves\_Fusionnes\_[Date].pdf}
        \item Contient tous les relevés dans l'ordre
    \end{itemize}
\end{itemize}

\subsection{Calcul automatique des moyennes et mentions}

\subsubsection{Calcul de la moyenne d'une UE}

Pour chaque Unité d'Enseignement (UE), la moyenne est calculée selon la formule :

\begin{tcolorbox}[colback=green!5!white,colframe=green!75!black,title=Formule de calcul]
$$\text{Moyenne UE} = \frac{\sum (\text{Note EC} \times \text{Poids EC})}{\sum \text{Poids EC}}$$

\textbf{Exemple :}

UE1 contient :
\begin{itemize}
    \item Mathématiques : note = 15/20, poids = 2
    \item Physique : note = 12/20, poids = 1.5
    \item Chimie : note = 14/20, poids = 2
\end{itemize}

$$\text{Moyenne UE1} = \frac{(15 \times 2) + (12 \times 1.5) + (14 \times 2)}{2 + 1.5 + 2} = \frac{30 + 18 + 28}{5.5} = \frac{76}{5.5} = 13.82$$
\end{tcolorbox}

\subsubsection{Calcul de la moyenne générale}

La moyenne générale du semestre est la moyenne pondérée des moyennes d'UE, pondérée par les crédits ECTS :

\begin{tcolorbox}[colback=green!5!white,colframe=green!75!black,title=Formule de calcul]
$$\text{Moyenne Générale} = \frac{\sum (\text{Moyenne UE} \times \text{Crédits UE})}{\sum \text{Crédits UE}}$$

\textbf{Exemple :}

\begin{itemize}
    \item UE1 : moyenne = 13.82, crédits = 30
    \item UE2 : moyenne = 11.5, crédits = 30
\end{itemize}

$$\text{Moyenne Générale} = \frac{(13.82 \times 30) + (11.5 \times 30)}{30 + 30} = \frac{414.6 + 345}{60} = \frac{759.6}{60} = 12.66$$
\end{tcolorbox}

\subsubsection{Attribution des mentions}

Les mentions sont attribuées automatiquement selon la moyenne générale :

\begin{table}[h]
\centering
\begin{tabular}{|l|c|l|}
\hline
\textbf{Mention} & \textbf{Moyenne} & \textbf{Couleur sur le relevé} \\
\hline
Très Bien & $\geq$ 16/20 & Vert foncé \\
\hline
Bien & 14 à 15.99 & Vert \\
\hline
Assez Bien & 12 à 13.99 & Bleu \\
\hline
Passable & 10 à 11.99 & Orange \\
\hline
Ajourné & < 10 & Rouge \\
\hline
\end{tabular}
\caption{Barème des mentions}
\end{table}

\subsection{Gestion des sessions de rattrapage}

Certains ECs peuvent avoir une session de rattrapage. Dans ce cas :

\begin{enumerate}[leftmargin=*]
    \item Lors de la configuration de l'EC dans \menu{Config des relevés}, cochez \champ{A une session}
    \item Dans le fichier Excel, créez deux colonnes :
    \begin{itemize}
        \item Une pour la note de session normale : \texttt{Mathematiques}
        \item Une pour la note de session de rattrapage : \texttt{Mathematiques\_Session}
    \end{itemize}
    \item Lors du mapping, associez les deux colonnes
    \item L'application prendra automatiquement la meilleure note entre les deux sessions
\end{enumerate}

\astuce{Si un étudiant n'a pas de note de rattrapage (cellule vide), seule la note de session normale sera utilisée.}

\subsection{Sélection d'étudiants spécifiques}

Vous pouvez générer des relevés pour certains étudiants uniquement :

\begin{enumerate}[leftmargin=*]
    \item Dans l'onglet \textbf{Sélection}, cochez \champ{Activer la sélection personnalisée}
    \item Trois options s'offrent à vous :

    \paragraph{Option 1 : Filtrer par niveau}
    \begin{itemize}
        \item Cochez les niveaux souhaités (L1, L2, L3, M1, M2...)
        \item Seuls les étudiants de ces niveaux seront traités
    \end{itemize}

    \paragraph{Option 2 : Filtrer par filière}
    \begin{itemize}
        \item Cochez les filières souhaitées
        \item Seuls les étudiants de ces filières seront traités
    \end{itemize}

    \paragraph{Option 3 : Sélection manuelle}
    \begin{itemize}
        \item Une liste complète des étudiants s'affiche
        \item Cochez individuellement les étudiants pour lesquels générer un relevé
        \item Utilisez la barre de recherche pour trouver rapidement un étudiant
    \end{itemize}

    \item Un compteur affiche le nombre d'étudiants sélectionnés
    \item Générez normalement avec \bouton{Générer les relevés}
\end{enumerate}

\begin{figure}[h]
    \centering
    % \includegraphics[width=\textwidth]{images/selection_etudiants.png}
    \fbox{\parbox{0.9\textwidth}{\centering [CAPTURE D'ÉCRAN : ONGLET SÉLECTION]\\ \vspace{0.3cm}
    ☑ Activer la sélection personnalisée\\
    Filtres : ☑ L1  ☐ L2  ☑ L3\\
    Liste avec cases à cocher par étudiant\\
    32 étudiants sélectionnés}}
    \caption{Sélection personnalisée d'étudiants}
\end{figure}

\subsection{Personnalisation du thème des relevés}

Vous pouvez personnaliser l'apparence des relevés de notes :

\begin{enumerate}[leftmargin=*]
    \item Dans le menu \menu{Config des relevés}, sélectionnez votre configuration
    \item Cliquez sur l'onglet \textbf{Thème}
    \item Configurez :
    \begin{itemize}
        \item \textbf{Couleurs} : Couleur primaire, secondaire, accent
        \item \textbf{Polices} : Famille de police, tailles
        \item \textbf{Mise en page} : Espacement, marges, disposition du contenu
    \end{itemize}
    \item Prévisualisez en temps réel les changements
    \item Enregistrez le thème
\end{enumerate}

Le thème sera automatiquement appliqué à tous les relevés générés avec cette configuration.

\newpage
